\chapter{Supervision temps réel : dashboard de télémétrie}
\label{chap:dashboard}

\section{Contexte}
% TODO : besoin d'un poste de supervision unique : état du véhicule, carte,
% lidar, latence réseau, messages V2V — visibles dans un simple navigateur.

\section{Problématique}
% TODO : exposer les données ROS 2 du véhicule à une interface web, en temps
% réel, à travers WiFi comme 5G, sans saturer le lien.

\section{Mission : Conception et développement du dashboard web}
\objectif
% TODO : spécifications des panneaux (carte SLAM, lidar, vitesse, latence
% WiFi/5G superposées, feed CAM/DENM).
\methodes
% TODO : architecture rosbridge (WebSocket :9090) + React/TypeScript +
% roslib ; nœud telemetry_bridge côté Jetson (ping, température CPU, mode
% wifi/5g) ; throttling des topics pour préserver la bande passante ;
% piège du ROS_DOMAIN_ID raconté en discussion.
\resultats
% TODO : captures d'écran du dashboard en fonctionnement (lidar live,
% courbes WiFi vs 5G superposées) ; usage pendant les démos.
\discussion
% TODO : choix React vs Foxglove justifié ; limites et évolutions (caméra,
% TF map->odom pour la pose exacte).

\section{Conclusion}
